\documentclass[../../../include/open-logic-section]{subfiles}

\begin{document}

\olfileid[fr]{sth}{cardinals}{classing}
\olsection{Finis, \usetoken{s}{enumerable} et \usetoken{s}{nonenumerable}}

Maintenant que nous connaissons les cardinaux, prenons
le temps d'en distinguer plusieurs sortes : les
cardinaux finis, !!{enumerable} et !!{nonenumerable}.

Nos deux premiers résultats entraînent que les
cardinaux finis sont exactement les ordinaux
finis, que nous avons définis comme nos
\emph{nombres naturels} dans la
\olref[z][infinity-again]{defnomega} :

\begin{prop}\ollabel{finitecardisoequal}
Si $n, m \in \omega$, alors $n = m$ si et
seulement si $\cardeq{n}{m}$.
\end{prop}

\begin{proof}
Le sens de gauche à droite est immédiat.
Pour l'autre sens, supposons
$\cardeq{n}{m}$ alors que $n \neq m$.
Par trichotomie, soit $n \in m$, soit
$m \in n$ ; supposons sans perte de
généralité que $n \in m$. Alors
$n \subsetneq m$ et il existe
!!a{bijection} $f \colon m \to n$ :
$m$ serait donc infini au sens de Dedekind,
contrairement à la
\olref[z][infinity-again]{naturalnumbersarentinfinite}.
\end{proof}

\begin{cor}\ollabel{naturalsarecardinals}
Si $n \in \omega$, alors $n$ est un cardinal.
\end{cor}

\begin{proof}
C'est immédiat.
\end{proof}
\noindent
Il s'ensuit aussi que plusieurs façons raisonnables
de qualifier un cardinal de «~fini~» ou d'«~infini~»
coïncident :
\begin{thm}\ollabel{generalinfinitycharacter}
Pour tout ensemble $A$, les propriétés suivantes
sont équivalentes :
\begin{enumerate}
	\item\ollabel{card:notinomega}
	$\card{A} \notin \omega$, c'est-à-dire que
	$A$ n'est pas fini ;\footnote{Note de l'édition
	française : l'original dit ici que $A$ n'est
	«~pas un nombre naturel~». Ce n'est pas
	équivalent à $\card{A}\notin\omega$ :
	un ensemble non ordinal peut être fini.
	Le sens voulu est que $A$ est infini.}
	\item\ollabel{card:omegaplus}
	$\omega \leq \card{A}$ ;
	\item\ollabel{card:infinite}
	$A$ est infini au sens de Dedekind.
\end{enumerate}
\end{thm}

\begin{proof}
Cela découle du
\olref[ord-arithmetic][using-addition]{ordinfinitycharacter},
du \olref[cardinals][cardsasords]{lem:CardinalsBehaveRight}
et du \olref{naturalsarecardinals}.
\end{proof}

Nous pouvons donc \emph{définir} rigoureusement
certaines notions employées assez librement dans la
\olref[sfr][][]{part} :

\begin{defn}\ollabel{defnfinite}
Nous disons que $A$ est \emph{fini} si et seulement si
$\card{A}$ est un nombre naturel, c'est-à-dire si
$\card{A} \in \omega$. Sinon, nous disons que
$A$ est \emph{infini}.
\end{defn}
\noindent
Notons toutefois que cette définition s'inscrit
dans le cadre de $\ZFC$. Il nous fallait en effet
l'axiome du bon ordre pour garantir que tout ensemble
a une cardinalité. Sans cet axiome, il peut exister
un ensemble qui n'est ni fini ni infini au sens de
Dedekind. Nous reviendrons sur cette question dans le
\olref[choice][]{chap}. Pour l'heure, nous
continuons d'employer l'axiome du bon ordre.

Passons des cardinaux finis aux cardinaux infinis.
Voici deux faits élémentaires :

\begin{cor}\ollabel{omegaisacardinal}
$\omega$ est le plus petit cardinal infini.
\end{cor}

\begin{proof}
$\omega$ est un cardinal : il est infini au sens
de Dedekind et, si $\cardeq{\omega}{n}$ pour
un $n \in \omega$, alors $n$ serait infini au
sens de Dedekind, contrairement à la
\olref[z][infinity-again]{naturalnumbersarentinfinite}.
Par définition, $\omega$ est donc le plus petit
cardinal infini.
\end{proof}

\begin{cor}
Tout cardinal infini est un ordinal limite.
\end{cor}

\begin{proof}
Soit $\alpha$ un ordinal successeur infini,
avec $\alpha = \beta \ordplus 1$.
D'après la \olref{finitecardisoequal},
$\beta$ est aussi infini ; donc
$\cardeq{\beta}{\beta \ordplus 1}$ par le
\olref[ord-arithmetic][using-addition]{ordinfinitycharacter}.
Ainsi,
$\card{\beta} = \card{\beta\ordplus 1}
= \card{\alpha}$ par le
\olref[cardinals][cardsasords]{lem:CardinalsBehaveRight},
et donc $\alpha \neq \card{\alpha}$.
\end{proof}

Dès la \olref[sfr][siz][enm-alt]{defn:enumerable},
nous avons signalé la différence entre ensembles
infinis !!{enumerable} et !!{nonenumerable}.
Cette définition conduit naturellement au résultat
suivant :

\begin{prop}
$A$ est !!{enumerable} si et seulement si
$\card{A} \leq \omega$, et $A$ est
!!{nonenumerable} si et seulement si
$\omega < \card{A}$.
\end{prop}

\begin{proof}
Par trichotomie, les deux affirmations sont
équivalentes ; il suffit donc de démontrer que
$A$ est !!{enumerable} si et seulement si
$\card{A} \leq \omega$. De droite à gauche :
si $\card{A} \leq \omega$, alors
$\cardle{A}{\omega}$ d'après le
\olref[cardinals][cardsasords]{lem:CardinalsBehaveRight}
et le \olref{omegaisacardinal}. De gauche à droite :
supposons $A$ !!{enumerable}. D'après la
\olref[sfr][siz][enm-alt]{defn:enumerable},
trois cas sont possibles :
\begin{enumerate}
	\item si $A = \emptyset$, alors
	$\card{A} = 0 \in \omega$, par le
	\olref{naturalsarecardinals} et le
	\olref[cardinals][cardsasords]{lem:CardinalsBehaveRight} ;
	\item si $\cardeq{n}{A}$, alors
	$\card{A} = n \in \omega$, par le
	\olref{naturalsarecardinals} et le
	\olref[cardinals][cardsasords]{lem:CardinalsBehaveRight} ;
	\item si $\cardeq{\omega}{A}$, alors
	$\card{A} = \omega$, par
	le \olref{omegaisacardinal}.
\end{enumerate}
Dans tous les cas, $\card{A} \leq \omega$.
\end{proof}
\noindent
$\omega$ occupe donc une place particulière.
Alors qu'il existe de nombreux ordinaux dénombrables :

\begin{cor}
$\omega$ est le seul cardinal infini !!{enumerable}.
\end{cor}

\begin{proof}
Soit $\cardfont{a}$ un cardinal infini
!!{enumerable}. Puisque $\cardfont{a}$
est infini, $\omega \leq \cardfont{a}$.
Puisque $\cardfont{a}$ est un cardinal
!!{enumerable},
$\cardfont{a} =
\card{\cardfont{a}} \leq \omega$.
Donc $\cardfont{a} = \omega$ par trichotomie.
\end{proof}

Il existe bien sûr une infinité de cardinaux.
Nous pourrions donc demander :
\emph{Combien y a-t-il de cardinaux ?}
Les résultats suivants suggèrent de revoir
la question.

\begin{prop}\ollabel{unioncardinalscardinal}
Si tout membre de $X$ est un cardinal,
alors $\bigcup X$ est un cardinal.
\end{prop}

\begin{proof}
Il est facile de vérifier que $\bigcup X$
est un ordinal. Soit
$\alpha \in \bigcup X$ un ordinal ; alors
$\alpha \in \cardfont{b} \in X$
pour un certain cardinal $\cardfont{b}$.
Comme $\cardfont{b}$ est un cardinal,
$\cardless{\alpha}{\cardfont{b}}$.
Puisque $\cardfont{b} \subseteq \bigcup X$,
nous avons
$\cardle{\cardfont{b}}{\bigcup X}$,
et donc
$\cardneq{\alpha}{\bigcup X}$.
En généralisant, $\bigcup X$ est un cardinal.
\end{proof}

\begin{thm}\ollabel{lem:NoLargestCardinal}
Il n'existe pas de plus grand cardinal.
\end{thm}

\begin{proof}
Pour tout cardinal $\cardfont{a}$, le théorème
de Cantor
(\olref[sfr][siz][car]{thm:cantor}) et le
\olref[cardinals][cardsasords]{lem:CardinalsExist}
impliquent que
$\cardfont{a} <
\card{\Pow{\cardfont{a}}}$.
\end{proof}

\begin{thm}
L'ensemble de tous les cardinaux n'existe pas.
\end{thm}

\begin{proof}
Par l'absurde, supposons que
$C = \Setabs{\cardfont{a}}
{\cardfont{a} \text{ est un cardinal}}$.
Alors $\bigcup C$ est un cardinal par la
\olref{unioncardinalscardinal}. D'après le
\olref{lem:NoLargestCardinal}, il existe donc
un cardinal $\cardfont{b} > \bigcup C$.
Par définition, $\cardfont{b} \in C$,
d'où $\cardfont{b} \subseteq \bigcup{C}$,
et donc $\cardfont{b} \leq \bigcup C$ :
contradiction.
\end{proof}

Comparez ce résultat au paradoxe de Russell
et au paradoxe de Burali-Forti.

\end{document}
